% One division, additively: the call places its wall, links the single pair the
% door falls on, and recurses. It never looks at the other crossings, because
% they were never there to remove. Here t = 0, the x axis, so the wall is vertical.
\begin{figure}[htbp]
  \centering
  \begin{subfigure}[b]{0.47\textwidth}\centering
    \begin{tikzpicture}[scale=0.82]
      \foreach \x in {0,1,2,3,4} \foreach \y in {0,1,2,3}
        \node[node vertex] (u\x\y) at (\x,-\y) {};
      \draw[door] (u21) -- (u31);
      \draw[wall] (2.5,0.45) -- (2.5,-0.75);
      \draw[wall] (2.5,-1.25) -- (2.5,-3.45);
      \node[font=\tiny, wallred, above] at (2.5,0.45) {wall at $j=2$};
      \node[font=\tiny, doorgreen, left] at (-0.2,-1) {door at $i{=}1$};
    \end{tikzpicture}
    \caption{The call links one pair and stops. Every other node is still
    isolated; no crossing was examined, let alone removed.}
    \label{fig:division-adds}
  \end{subfigure}\hfill
  \begin{subfigure}[b]{0.47\textwidth}\centering
    \begin{tikzpicture}[scale=0.82]
      \foreach \x in {0,1,2,3,4} \foreach \y in {0,1,2,3}
        \node[node vertex] (v\x\y) at (\x,-\y) {};
      % left half, carved by the first recursive call
      \foreach \a/\b in {v00/v10, v10/v20, v00/v01, v01/v02, v02/v03,
                         v20/v21, v11/v21, v02/v12, v12/v22,
                         v03/v13, v13/v23}
        \draw[black!55] (\a) -- (\b);
      % right half, carved by the second
      \foreach \a/\b in {v30/v40, v30/v31, v31/v32, v32/v42, v42/v43,
                         v32/v33, v40/v41}
        \draw[black!55] (\a) -- (\b);
      \draw[door] (v21) -- (v31);
      \draw[wall] (2.5,0.45) -- (2.5,-0.75);
      \draw[wall] (2.5,-1.25) -- (2.5,-3.45);
    \end{tikzpicture}
    \caption{After both recursive calls return. Each half is a tree, and the
    door is still the only edge between them.}
    \label{fig:division-recursed}
  \end{subfigure}
  \caption{One division, with $t=0$ --- the coin named the $x$ axis, so the wall
  is vertical. The call draws the wall position $j$ on that axis and the door
  position $i$ on the other, links the two cells the door falls on, and recurses
  on the two halves. Counting edges: $11$ on the left, $7$ on the right, one
  door, total $19=20-1$ on $20$ cells --- the tally of
  Proposition~\ref{prop:generation-is-a-tree}.}
  \label{fig:one-division-step}
\end{figure}
